\documentclass{article}
\usepackage[OT1]{fontenc} % Match official pdfLaTeX Times encoding under XeTeX as well.
\usepackage{microtype}
\usepackage{graphicx}
\usepackage{booktabs}
\usepackage{amsmath}
\usepackage{array}
\usepackage{tabularx}
\usepackage{xurl}
\usepackage[hypertexnames=false]{hyperref}
\usepackage{mlsys2025}
\hypersetup{pdfauthor={Anonymous Authors},pdftitle={Supplement: Execution-Path Qualification and Realized Costs in Speculative Decoding on Consumer Systems},pdfsubject={MLSys 2027 separate appendix v0.8.6}}
% Generated from repository evidence by scripts/build_assets.py.
\newcommand{\MRepSpeed}{0.842}
\newcommand{\MRepN}{24}
\newcommand{\MRepPairs}{72}
\newcommand{\MRepExact}{72}
\newcommand{\MRepFaster}{1}
\newcommand{\MRepLo}{0.809}
\newcommand{\MRepHi}{0.877}
\newcommand{\MHEspeed}{1.387}
\newcommand{\MHEprompts}{30}
\newcommand{\MHEpairs}{90}
\newcommand{\MHEexact}{90}
\newcommand{\MHEfaster}{30}
\newcommand{\MHElo}{1.351}
\newcommand{\MHEhi}{1.422}
\newcommand{\MHEenergy}{1.013}
\newcommand{\MHEenergyLo}{0.990}
\newcommand{\MHEenergyHi}{1.034}
\newcommand{\RTXcpuSpeed}{0.404}
\newcommand{\RTXcpuLo}{0.390}
\newcommand{\RTXcpuHi}{0.418}
\newcommand{\RTXgpuSpeed}{0.650}
\newcommand{\RTXgpuLo}{0.629}
\newcommand{\RTXgpuHi}{0.673}
\newcommand{\RTXexact}{90}
\newcommand{\BSpeed}{0.3087}
\newcommand{\BLo}{0.2927}
\newcommand{\BHi}{0.3246}
\newcommand{\BvShort}{1.983}
\newcommand{\BvLong}{1.992}
\newcommand{\BExact}{60}
\newcommand{\BExactPrompts}{20}
\newcommand{\BFastExact}{45}
\newcommand{\RTXrawZero}{1.009}
\newcommand{\RTXserZero}{2.431}
\newcommand{\RTXrawDepth}{0.983}
\newcommand{\RTXserDepth}{2.137}
\newcommand{\HExact}{90}
\newcommand{\HcpuSpeed}{0.734}
\newcommand{\HcpuLo}{0.693}
\newcommand{\HcpuHi}{0.771}
\newcommand{\HcpuC}{2.4631}
\newcommand{\HcpuDraft}{12.63}
\newcommand{\HcpuTarget}{87.34}
\newcommand{\HcpuResidual}{0.04}
\newcommand{\HcpuD}{0.4304}
\newcommand{\HcpuR}{2.9767}
\newcommand{\HcpuX}{0.00124}
\newcommand{\HcpuAggregate}{0.7227}
\newcommand{\HnpuSpeed}{0.352}
\newcommand{\HnpuLo}{0.333}
\newcommand{\HnpuHi}{0.369}
\newcommand{\HnpuC}{2.4584}
\newcommand{\HnpuDraft}{58.35}
\newcommand{\HnpuTarget}{41.63}
\newcommand{\HnpuResidual}{0.02}
\newcommand{\HnpuD}{4.1251}
\newcommand{\HnpuR}{2.9435}
\newcommand{\HnpuX}{0.00130}
\newcommand{\HnpuAggregate}{0.3477}

% Generated from frozen runtime evidence by runtime_values.py.
\newcommand{\RTGuard}{1.3114}
\newcommand{\RTGuardLo}{1.2487}
\newcommand{\RTGuardHi}{1.3787}
\newcommand{\RTAlways}{1.3469}
\newcommand{\RTAlwaysLo}{1.2867}
\newcommand{\RTAlwaysHi}{1.4088}
\newcommand{\RTOracle}{1.0278}
\newcommand{\RTOracleLo}{1.0137}
\newcommand{\RTOracleHi}{1.0446}
\newcommand{\RTFallback}{30.6}
\newcommand{\RTK}{2}
\newcommand{\RTCap}{128}
\newcommand{\RTRounds}{8}
\newcommand{\RTMargin}{0}
\newcommand{\RTCalibration}{40}
\newcommand{\RTQualification}{16}
\newcommand{\RTPrompts}{60}
\newcommand{\RTQualPairs}{48}
\newcommand{\RTReps}{3}
\newcommand{\RTObservations}{540}
\newcommand{\RTExact}{180}
\newcommand{\RTBootstrap}{10,000}
\newcommand{\RTSeed}{20260924}
\newcommand{\RTSecondary}{1.0271}
\newcommand{\RTSecondaryLo}{1.0130}
\newcommand{\RTSecondaryHi}{1.0436}
\newcommand{\RTGuardFaster}{24}
\newcommand{\RTAlwaysFaster}{36}
\newcommand{\RTTies}{0}

% Generated by scripts/direct_id_values.py from frozen formal JSON; do not edit.
\newcommand{\DIRTXBatchedPairs}{75}
\newcommand{\DIRTXBatchedPrompts}{25}
\newcommand{\DIRTXBatchedStatus}{BLOCKED}
\newcommand{\DIRTXSerializedPairs}{90}
\newcommand{\DIRTXSerializedPrompts}{30}
\newcommand{\DIRTXSerializedStatus}{QUALIFIED}
\newcommand{\DIBFastPairs}{60}
\newcommand{\DIBFastPrompts}{20}
\newcommand{\DIBFastStatus}{BLOCKED}
\newcommand{\DIBSerializedPairs}{78}
\newcommand{\DIBSerializedPrompts}{26}
\newcommand{\DIBSerializedStatus}{BLOCKED}

% New campaign projection; historical macros remain unchanged.
\newcommand{\FUCost}{0.646}
\newcommand{\FUCostLo}{0.630}
\newcommand{\FUCostHi}{0.660}
\newcommand{\FUForkScores}{72}
\newcommand{\FUForkImmediate}{0}
\newcommand{\FUForkLater}{18}
\newcommand{\FUForkLaterPrompts}{6}
\newcommand{\FUForkEvents}{195}

% Post-study existing-record derived values; no original primary macro altered.
\newcommand{\DerCPUZero}{0.8275}
\newcommand{\DerNPUZero}{0.8352}
\newcommand{\DerAlwaysOracle}{1.0007}
\newcommand{\DerTargetFaster}{2}
\newcommand{\DerScoreAbsMedian}{0.0695}
\newcommand{\DerScoreAbsMax}{2.3503}
\newcommand{\DerMarginMin}{0.2291}
\newcommand{\DerMarginAltMin}{0.2003}
\newcommand{\DerRTXTargetMs}{1907.9}
\newcommand{\DerRTXSerialMs}{2923.3}

\mlsystitlerunning{Supplement: Execution-Path Qualification in Speculative Decoding}
\begin{document}
\twocolumn[{
\mlsystitle{Supplement: Execution-Path Qualification and Realized Costs in Speculative Decoding on Consumer Systems}
\begin{mlsysauthorlist}\mlsysauthor{Anonymous Authors}{anon}\end{mlsysauthorlist}
\mlsysaffiliation{anon}{Anonymous Institution}
\mlsyscorrespondingauthor{Anonymous Author}{anon.email@domain.com}
\vskip 0.3in
}]
\printAffiliationsAndNotice{}
\appendix
\renewcommand{\thetable}{A\arabic{table}}
\section{Supplement Scope}
This separate appendix accompanies the measured hardware, mechanism, and runtime studies. It supplies statistical, measurement-boundary, and provenance details. Main B570 H1 maps to Section J and Tables A6--A7; new forks to Section L and the post-study tables; RTX prediction to Section M and its input audit; RTX cost to Section N; M4 to Section H and the post-study fixed-mode comparison; 358H to Section D and the conditional accounting. Section O identifies exact packages and verifier entries. Identifying repository metadata and internal source paths remain outside the anonymous PDF; the nonidentifying B570 prediction-file content digest appears below.

\section{Frozen Estimators}
Table~\ref{tab:statistics} (A1) preserves historical hardware estimators only; M4 runtime is in Section H, new RTX cost in Section N, and new post-study sensitivity in Section~\ref{app:derived}. The primary sample unit is the prompt. Repeated observations contribute to its summary, rather than being treated as independent prompts. Pair ratios divide matched baseline decode time by speculative decode time. Historical M4, RTX, and B570 counts are text-path agreement from parsed CLI output: M4 HumanEval includes the prompt in its parsed text; B570 strips and normalizes the prompt before re-tokenization. Neither is direct generated-ID equality. The 358H runner and prospective M4 study record generated IDs directly. A geometric mean across prompt summaries is used for headline speedup. B570 timing remains diagnostic regardless of its estimator.

\begin{table*}[t]
\caption{Frozen statistical policies. GM = geometric mean. Indices below are zero-based order statistics of sorted bootstrap draws; NumPy quantiles use the original analyzer's convention.}
\label{tab:statistics}
\centering
\begin{tabular}{>{\raggedright\arraybackslash}p{1.1in}>{\raggedright\arraybackslash}p{1.1in}>{\raggedright\arraybackslash}p{1.0in}>{\raggedright\arraybackslash}p{.9in}>{\raggedright\arraybackslash}p{1.65in}}
\toprule
Campaign & Within prompt & Resamples & Seed & CI endpoints / generator \\
\midrule
M4 24Q & GM of 3 ratios & 10,000 & 20260907 & [249,9749]; Python Random \\
M4 HumanEval & GM of 3 ratios & 100,000 & 20260910 & [2500,97499]; Python Random \\
RTX HumanEval & GM of 3 ratios & 10,000 & 20260910 & [249,9749]; Python Random \\
B570 fallback & GM of 3 ratios & 10,000 & 20260910 & [250,9750]; Python Random \\
358H Phase 3 & Median of 3 ratios & 10,000 & 20260919 & Quantiles .025/.975; NumPy default\_rng \\
\bottomrule
\end{tabular}
\end{table*}
The B570 finalizer builds its prompt vector in the insertion order of the archived pair records. Sorting by task ID changes the finite fixed-seed bootstrap endpoints even though it leaves the underlying estimator unchanged. Exact replay retains that original order. The M4 HumanEval lower endpoint uses a different order statistic from M4 24Q and RTX; those frozen conventions are retained. None of these prompt intervals captures the full distribution of possible thermal, session or runtime states.

\section{Energy Boundaries and Allowed Claims}
% Evidence: M4-HE/SECONDARY; B570-SERIAL; RTX reports; 358H Phase3.
\begin{table*}[t]
\caption{Sensor boundaries and allowed interpretations. No cells are combined into whole-system energy.}
\label{tab:energy}
\centering
\begin{tabular}{>{\raggedright\arraybackslash}p{1.35in}>{\raggedright\arraybackslash}p{2.1in}>{\raggedright\arraybackslash}p{2.45in}}
\toprule
Campaign & Measurement boundary & Allowed statement \\
\midrule
M4 HumanEval & Silemetry CPU + GPU + ANE package aggregate; separate DRAM excluded & Fixed-$k=2$ package-energy interval includes parity; registered threshold variant secondary \\
RTX 5070 Ti Laptop & NVIDIA board / CPU package are distinct domains; follow-up process window includes loading & No formal serialized HumanEval energy comparison \\
B570 fallback & Arc adapter board only; no CPU package or wall power & Diagnostic board efficiency only; path not globally exact \\
358H Phase 3 & No qualified end-to-end energy boundary & No end-to-end CPU claim; NPU energy not admissible \\
\bottomrule
\end{tabular}
\end{table*}

Ratios above one favor speculation within the named boundary. For M4 fixed-$k=2$, the package ratio is \MHEenergy\ [\MHEenergyLo, \MHEenergyHi], so a package-energy improvement is not established. The registered $p_{\min}=0.9$ secondary condition has a package-efficiency ratio 1.041 [1.018, 1.065]. This is not a replacement for the primary condition and does not establish a material additional speed gain. B570 serialized board efficiency is 0.302 [0.285, 0.320]; correctness failure restricts it to diagnosis. No between-platform energy ranking is valid.

\section{Round Accounting and State Semantics}
The 358H target keeps its low-level context across rounds. Candidate evaluation, invalid-suffix rollback when needed, and correction or bonus decoding belong to the complete target-side timing region. Every OpenVINO proposal call receives the entire committed prefix. Its synchronous timing combines reconstruction and proposal; neither term can be isolated from this experiment without changing the measured execution.

The CPU path records 11,520 committed output tokens over 4,677 rounds; the NPU path records 11,520 over 4,686 rounds. Their directly measured $C$ values are \HcpuC\ and \HnpuC. Using the arithmetic baseline time per output token and arithmetic component totals yields CPU $(D,R,X)=(\HcpuD,\HcpuR,\HcpuX)$ and NPU $(\HnpuD,\HnpuR,\HnpuX)$. The resulting ratios \HcpuAggregate\ and \HnpuAggregate\ match ratios of total times, not the prompt-GM headline estimator.

Historical M4/RTX accepted-draft counters exclude correction or bonus output and have output-boundary effects. They are not substituted for the directly recorded 358H committed-progress metric. Likewise, target-forward PP2/PP1 ratios are never presented as complete realized target-round cost.

\section{Deviations and Provenance Boundaries}
The M4 corrected synthetic ceiling was preregistered but not executed. Earlier synthetic acceptance diagnostics are not a substitute. The representative and separately preregistered code workloads stay distinct. Early MLX precision/correctness failures are qualification history, not failures of the final Metal Q8\_0 path.

The B570 formal v2 correctness failure triggered its stop rule. A later six-prompt serialized probe passed 18/18 text-path comparisons, but full-workload serialization yielded only 60/90 under the same comparator. Earlier remote-interaction-contaminated verification observations are excluded from the plotted final serialized cells. A legacy ``exact'' field label in historical artifacts does not establish workload-global equality.

Windows follow-up derived tables and numerical captures support the manuscript, but historical raw per-run provenance is weaker. The CUDA score capture came from a separate same-commit instrumented build and does not replace formal measurements. The 358H source reconstruction is strongly supported by retained evidence but lacks proof of exact historical upstream identity (the internal \texttt{STRONGLY\_SUPPORTED} label). Its completed Phase-3 records supersede the blocked-preflight empty scaffold. These are retained qualifications, not unresolved invitations to alter frozen experiments.

The historical RTX HumanEval/106 CUDA reproducer selects token 2790 under sequential/serialized execution and 2629 under batching in three of three runs at the same location. A separate same-commit diagnostic build captures the score-margin change from $+0.080948$ to $-0.126457$. This supports stable batch-shape-dependent greedy argmax divergence in that historical configuration; it does not establish a specific kernel, universal CUDA defect, or a shared RTX/B570 root cause.
\section{Verification Context Lookup}
Table~\ref{tab:verificationlookup} retains all historical context cells behind Panel A of the main verification figure. Values are target-forward ratios, not complete target-side round cost or integrated speedup. Context ranges do not represent uncertainty intervals. RTX denotes RTX 5070 Ti Laptop throughout this appendix.
\begin{table*}[t]
\caption{Verification-path lookup. Each row is a measured context cell; ratios are shown to six decimal places from the frozen source. Qualification and equality labels refer to the historical text-path comparator for the associated execution path; they do not certify generated-ID equality. M4 contexts and B570 contexts are prefix lengths; RTX uses the reported context depth. The stacks and context sets are not matched. The original rounded ranges are M4 Metal 1.02--1.08, RTX batched 0.98--1.01, RTX serialized 2.14--2.43, B570 fast 1.24--1.64, and B570 serialized 1.98--1.99; these span contexts, not confidence intervals.}
\label{tab:verificationlookup}
\centering
\begin{tabular}{llrl}
\toprule
Path & Context (tokens) & $V_2^{\mathrm{micro}}$ & Qualification \\
\midrule
M4 Metal & 512 & 1.022182 & Qualified shared path \\
M4 Metal & 2048 & 1.076126 & Qualified shared path \\
M4 Metal & 8192 & 1.032264 & Qualified shared path \\
RTX batched & 0 & 1.009030 & Correctness-blocked; diagnostic \\
RTX batched & 256 & 0.983148 & Correctness-blocked; diagnostic \\
RTX serialized & 0 & 2.430996 & Exact secondary control \\
RTX serialized & 256 & 2.137011 & Exact secondary control \\
B570 fast & 512 & 1.242576 & 45/90 exact; diagnostic, inadmissible \\
B570 fast & 2048 & 1.637295 & 45/90 exact; diagnostic, inadmissible \\
B570 serialized & 512 & 1.982668 & 60/90 exact; diagnostic, inadmissible \\
B570 serialized & 2048 & 1.991707 & 60/90 exact; diagnostic, inadmissible \\
\bottomrule
\end{tabular}
\end{table*}


\section{Evidence Inventory}
Table~\ref{tab:regimes} (A4) collects only the original case inventory and evidence classes; follow-up cohorts and endpoints are documented in Sections K--N. It is not an inventory of all newer prediction studies. These are case descriptions, not a hardware ranking, population-wide failure rate, or exhaustive taxonomy. M4's separate-session studies establish opposite observed outcomes; neither diagnosed mechanism is assigned to that contrast.
% Immutable table body from merged scientific style PR #19.
\begin{table*}[t]
\caption{Evidence inventory with row-level output endpoints. P: primary; S: secondary; M: mechanism; D: diagnostic. Historical M4/RTX/B570 rows report text-path agreement; 358H reports direct-ID sequence equality. Both B570 paths remain inadmissible diagnostics. The RTX serialized controls each have 90/90 historical text-path agreement. The 358H P/M entries separate primary outcomes from mechanism-supported decomposition. None establishes functional code correctness; annotations are descriptive and non-exhaustive.}
\label{tab:regimes}
\centering
\begingroup
% AUTO-GENERATED FROM FROZEN STYLE TABLE; OUTPUT ENDPOINT LABELS ONLY.
\setlength{\tabcolsep}{4pt}
\renewcommand{\arraystretch}{1.15}
\begin{tabularx}{\linewidth}{@{}p{1.38in}p{1.25in}p{.95in}p{.35in}X@{}}
\toprule
Case & Evidence endpoint / gate & Outcome & Class & Case annotation \\
\midrule
M4 representative 24Q & Text-path 72/72 & Negative & P & Workload-conditioned \\
M4 HumanEval30 & Text-path 90/90 & Profitable & P & Workload-conditioned \\
RTX batched & Text-path; BLOCKED & Not admissible & M & Verification-constrained \\
RTX serialized controls & Text-path 90/90 each & Negative & S & Verification-constrained \\
B570 fast & Text-path 45/90 & Not admissible & D & Correctness-blocked \\
B570 serialized & Text-path 60/90 & Not admissible & D & Correctness-blocked \\
358H CPU & Direct-ID 90/90 & Negative & P/M & Target-side dominated\newline 87.34\% target side \\
358H NPU & Direct-ID 90/90 & Negative & P/M & State-interface bound\newline 58.35\% replay + proposal \\
\bottomrule
\end{tabularx}

\endgroup
\end{table*}

\section{Runtime Study Reproducibility}
The M4 runtime estimator first takes each condition's median of \RTReps\ repetitions, then forms a prompt latency ratio and its geometric mean across \RTPrompts\ prompts. The offline oracle uses the smaller target-only or Always-Spec median for each prompt. Bootstrap intervals use \RTBootstrap\ draws with Python Random seed \RTSeed, resampling prompts in the frozen held-out order. Percentile endpoints linearly interpolate the sorted bootstrap draws at quantiles .025 and .975. A single random stream computes Target/Guard, Target/Always-Spec, and Guard/oracle intervals in that order. This preserves the frozen finite-bootstrap calculation; it does not replace the hardware estimators in Table~\ref{tab:statistics}. Intervals condition on the selected policy and observed prompt medians; they do not incorporate recalibration or between-session variation. The oracle selects between two measured fixed conditions, not all possible switching policies, so it is not a universal latency lower bound.

Calibration, qualification, and held-out IDs are disjoint. Selection ranks tasks with no recorded prior M4 exposure in the frozen inventory by a seeded hash within each of HumanEval and Spec-Bench; each split is balanced across these two families. The registered grid has nine candidates: first decision round $E\in\{2,4,8\}$ and margin $m\in\{0,0.05,0.10\}$. Calibration selects the exact candidate with the lowest geometric-mean latency ratio to the offline oracle. Exact ties break by mean measured guard-compute overhead, smaller exploration count, larger margin, and lexicographic candidate ID. The frozen winner is \texttt{E8\_M000}; held-out results never select or retune it.

The portable M4 archive contains original observations with generated-token IDs, complete-attempt manifests, qualification and calibration records, frozen source/build information, and analysis. An independent offline verifier reproduces policy selection, qualification and formal direct-ID equality, prompt medians, ratios, intervals, and fallback fraction without executing the models. No runtime energy result is claimed.

The reviewer-requested Guard/Always-Spec contrast is post-study paired secondary analysis, not the registered primary endpoint. Aggregate results were already known; the committed analysis plan is a procedural timestamp, not preregistration or blinding. For each frozen held-out prompt $i$, form $r_i=\operatorname{median}(T_{\mathrm{Guard},i})/\operatorname{median}(T_{\mathrm{Always},i})$ over its \RTReps\ repetitions, then compute $\exp(\RTPrompts^{-1}\sum_i\log r_i)$. Values above one mean Guard is slower. A fresh standalone Python Random stream with seed \RTSeed\ resamples \RTPrompts\ paired prompts per draw in frozen held-out order, using \RTBootstrap\ draws and the same linear percentile interpolation; it does not continue the primary stream. The ratio is \RTSecondary\ [\RTSecondaryLo, \RTSecondaryHi], with the interval above parity. Descriptive prompt counts are \RTGuardFaster\ Guard-faster, \RTAlwaysFaster\ Always-Spec-faster, and \RTTies\ exact ties. Both new implementations reconstruct raw medians from the same hash-verified archive and match the primary medians exactly; an independent calculation reproduces both finite-bootstrap endpoints. No new measurement was made. The comparison and its analysis choices were fixed for this secondary replay after aggregate results were already known; they were not blinded or prespecified in the original experiment. No choices were changed during that replay to improve the result. The v0.8.6 sensitivity below is a separate, explicitly post-study choice. The interval excludes session and recalibration variation; source hashes and all prompt ratios remain in the private evidence ledger.

A metadata-recording defect requires explicit provenance recovery. Loop-variable shadowing caused the frozen helper to return the build-manifest path rather than the calibration-freeze path. Qualification and formal-start records therefore contain the build-manifest digest in the intended freeze-hash field. A versioned recovery analyzer verified the actual freeze against the clean formal-start commit, qualification, frozen source/model/calibration/policy hashes, and complete attempts before applying the original formulas. The archive retains the original error markers and pipeline \texttt{STOPPED} state, meaning automated finalization stopped at the metadata exception. The scientific policy itself was not retuned. Its actual calibration-freeze digest is \nolinkurl{3c179a35f046b57edb4ddc8772a1aba609b496878efc2b7849a66947c8b7830b}; the misrecorded build-manifest digest is \nolinkurl{21592be38c0c477f63cea819b96b5b12cfea4529c131f6914b761a06ae5e518b}. Recovery compares the actual freeze and qualification bytes to the clean formal-start commit, rechecks source/model/calibration/policy identities and complete attempts, then applies the unchanged estimator. It validates retained links and reports recorded ordering; it does not supply trusted external time. This is a metadata-recording defect, not a measurement mutation.

\begin{table*}[t]
\caption{Fresh direct-ID qualification studies. Each path must achieve 90/90 complete generated-ID sequence equality; a wholly equal prompt matches in all three repetitions. QUALIFIED/BLOCKED is the frozen gate result. These studies have no matching fresh latency endpoint.}
\label{tab:directidqualification}
\centering
\begin{tabular}{llllll}
\toprule
Platform & Path & Direct-ID pairs & Equal prompts & Gate & Evidence endpoint \\
\midrule
RTX 5070 Ti Laptop & Batched & \DIRTXBatchedPairs/90 & \DIRTXBatchedPrompts/30 & \DIRTXBatchedStatus & Direct-ID sequences \\
RTX 5070 Ti Laptop & Serialized & \DIRTXSerializedPairs/90 & \DIRTXSerializedPrompts/30 & \DIRTXSerializedStatus & Direct-ID sequences \\
B570 & Fast & \DIBFastPairs/90 & \DIBFastPrompts/30 & \DIBFastStatus & Direct-ID sequences \\
B570 & Serialized & \DIBSerializedPairs/90 & \DIBSerializedPrompts/30 & \DIBSerializedStatus & Direct-ID sequences \\
\bottomrule
\end{tabular}
\end{table*}

The 358H diagnostic is pinned separately from the hardware campaign. Its compact evidence records reproducible target-only same-prefix divergence and proposal playback without live NPU inference, but does not establish a specific rollback-boundary bug. It explicitly records that no formal runtime performance data exist and that failed qualification prevented formal runtime performance evaluation. Internal evidence records retain the exact source commit and file hashes; the local raw diagnostic traces are not included in that compact record. No 358H runtime speedup is inferred from these diagnostics.
\section{Prospective Direct-ID Qualification}
Table~\ref{tab:directidqualification} reports fresh direct-ID qualification separately from the historical text-path lookup and timing inventory. Both studies freeze the same HumanEval30 IDs, Qwen3-8B and Qwen3-0.6B Q8\_0 models, greedy $k=2$, three repetitions, cap 128, and a 90/90 complete direct-ID sequence-equality rule. Prompt IDs and uncommitted proposals are excluded; an earlier generated terminal EOG is included. No mismatch is discarded and no subset qualifies. Historical path-control outcomes informed the replications; prospective freezing does not imply blinded discovery. Each study completes 270 executions; independent offline implementations reproduce all 180 comparisons, including first divergences. Operational elapsed fields do not establish formal performance.

The RTX original embedded-server CLI exposes native committed-token debug events without a source patch. It uses context 4096, reference target micro-batch one, candidate micro-batches two/one, driver 616.92, and full CUDA0 offload. The historical RTX driver was 616.56; this is not an exact replay of historical host state. The incomplete RTX v1 transport attempt is preserved and supplies no v2 formal pair; completed v2 uses lossless byte transport, a new prospective freeze, and a complete fresh schedule.

B570 uses source instrumentation to record committed IDs in sidecars, context 1024, and a separately frozen completion reference. Its endpoint selects the first 128 committed IDs or the inclusive earlier EOG; raw CLI overshoot is retained. Vulkan0 uses driver 32.0.101.9033 with full target 37/37 and draft 29/29 layer offload. The initial Task Scheduler launch failed during CP936/UTF-8 environment-metadata decoding before inference, with zero formal executions. The first data-producing attempt directly ran the unchanged frozen runner: no semantic, code, binary, model, driver, or parameter change occurred. This deployment deviation is preserved in provenance.

The runtime/front-end, reference, and context differences preclude a causal hardware comparison. Neither study estimates population failure rates or functional code correctness. Repeat stability holds for every prompt in each condition, but stable disagreement remains disagreement. Historical B570 fast 45/90 and serialized 60/90 text-path counts and historical RTX controls retain their original comparators. New qualification does not retrospectively authorize old timing as strict-ID performance. Portable validators hash included evidence and check recorded identities of omitted external weights/binaries; they do not rehash omitted assets locally.

\section{B570 Target-Partition Prediction}
The B570 Vulkan follow-up uses Qwen3-8B/0.6B Q8\_0, greedy $k=2$, context 1024 and 128 committed IDs or earlier inclusive EOG. Its target-only harness matches the UB128 reference's input IDs, sampler and APIs; H1 changes only target $n_{\mathrm{ubatch}}$ to one. The harness reproduces 90/90 reference sequences before comparison with live serialized speculation.

\paragraph{Selection and freeze.}
The audit covers 8,706 retained file records and archive members; any B570 association is conservatively excluded, even if execution is uncertain. This excludes 31 of 164 tasks. The 24 reserved IDs have the smallest SHA-256 digest of a fixed study prefix plus each eligible task ID; outcomes play no role. Discovery uses exposed HumanEval/89, 101, 104, 162, 3, 10, 68, 70, 19, 20, 28, 32. The frozen tree selects H1 at 36/36; only failure permits H2 fixed-proposal replay (also 36/36), and failure of both stops validation. H1 passes; H2 is not built. Predictor, endpoint, reserved IDs, binaries and criterion are committed before reserved inference.

\paragraph{Prospective validation.}
The 24 reserved IDs are HumanEval/58, 156, 134, 42, 122, 114, 93, 115, 151, 63, 125, 133, 43, 139, 77, 35, 27, 34, 145, 82, 154, 78, 96, and 138. Target reference and target-only predictor run three times each (144 executions) with within-condition repeat stability. All 72 predicted complete ID arrays and their digest are committed before the first of 72 live serialized outcomes. Primary mechanism validation requires 72/72 predictor-to-live complete-ID agreement; no partial threshold or confidence interval is substituted. That agreement is 72/72. Twenty prompts are predicted reference-equal; four are predicted divergent, and all 24 classifications match. Among the 12 divergent repetition pairs, the first-divergence index matches 12/12. Those four prompts and zero-based indices (reference token $\rightarrow$ serialized token) are HumanEval/134: 88 (914$\rightarrow$1946); /42: 73 (432$\rightarrow$1817); /145: 103 (1184$\rightarrow$686); and /138: 93 (14374$\rightarrow$9064). Reference-to-live direct-ID qualification is 60/72 and BLOCKED on this reserved set.

An independent validator passes 5,681 input, output, hash, endpoint and chronology checks; extracted archive validation passes. No prediction edits, selective reruns, parameter changes, performance or energy experiments occur after the relevant freezes. No-known-exposure is limited to audited records. The result establishes tested partition sufficiency, not a unique kernel, universal rule, CUDA/B570 shared cause or production certificate.
\paragraph{Reviewer-visible commitment record.}
% AUTO-GENERATED FROM VERIFIED PHASE-1 LEDGER.
% DO NOT EDIT RESULT VALUES BY HAND.
The comparison uses complete generated-token-ID arrays; prompt IDs are excluded.
At most 128 committed generated IDs are compared, including an earlier terminal end-of-generation (EOG) ID when it ends generation.
No mismatch is discarded. Equality requires the complete sequences, including lengths, to match; first-divergence indices are zero-based.

\noindent\textbf{Content identity.} The 72 frozen complete prediction arrays are bound to the following prediction-file content digest (SHA-256; two consecutive 32-character lines):
\begin{center}\small\texttt{0201ccd1f67492510d34be567aa83f1c}\\\texttt{38d3bd1f162a970c05d3cea9d3957826}\end{center}
The digest identifies the frozen contents; the temporal ordering below comes from the recorded study chronology rather than from the digest alone.

\begin{center}
\small
\begin{tabular}{@{}ll@{}}
\toprule
Recorded event & Time relative to H1 freeze \\
\midrule
H1 freeze & $T_0$ \\
72 prediction arrays finalized & $T_0 + 1722.38\,\mathrm{s}$ \\
First live reserved outcome & $T_0 + 1813.75\,\mathrm{s}$ \\
\bottomrule
\end{tabular}
\end{center}
The recorded study chronology fixes the prediction file before the first live serialized reserved outcome, with a 91.37\,s interval between prediction completion and live start. Identifying repository metadata are omitted for double-blind review.

\begin{table}[t]
\caption{Reference trajectory agreement versus frozen UB1 predictor agreement on the same 24 reserved prompts $\times$ 3 repetitions. Predictor/live complete-sequence equality is the primary mechanism-validation endpoint; reference/live equality is an already-observed comparison. The 60/72 reference result remains BLOCKED under its qualification contract.}
\label{tab:b570-commitment-summary}
\centering\small
\begin{tabular}{@{}lr@{}}
\toprule
Comparator to live serialized output & Complete equality \\
\midrule
Frozen UB128/reference trajectory & 60/72 \\
Frozen UB1 target-only predictor & 72/72 \\
\bottomrule
\end{tabular}
\end{table}

\noindent\textbf{Informative reference-divergent prompts.} These cases are reference-divergent but predictor/live-identical.
\begin{center}\small
\begin{tabular}{@{}lrrr@{}}
\toprule
Task & First index & Ref $\rightarrow$ Pred/Live ID & Pred=Live \\
\midrule
HE/134 & 88 & 914 $\rightarrow$ 1946 & 3/3 \\
HE/42 & 73 & 432 $\rightarrow$ 1817 & 3/3 \\
HE/145 & 103 & 1184 $\rightarrow$ 686 & 3/3 \\
HE/138 & 93 & 14374 $\rightarrow$ 9064 & 3/3 \\
\bottomrule
\end{tabular}
\end{center}
Across these 4 prompts, the predicted and observed first-divergence positions agree in 12/12 repetition pairs.

\makeatletter\global\setlength{\@dblfptop}{0pt}\makeatother
\begin{table*}[!t]
\caption{B570 reserved commitment ledger. All 24 prompts have three repeat-stable 128-ID outputs. Pred=Live and Ref=Live count complete generated-ID equality across three runs. Each displayed digest is a 12-hex prefix of the full SHA-256 shared by predictor and live output; complete arrays and hashes are in the anonymous sequence supplement. The first reference-to-live divergence index is zero-based; -- means reference equality.}
\label{tab:b570-commitment-ledger}
\centering\small
\begin{tabular}{@{}lccccr@{}}
\toprule
Task & Live len & Pred=Live & Ref=Live & Pred/Live digest & Ref $\rightarrow$ Live first divergence \\
\midrule
HE/58 & 128 & 3/3 & 3/3 & \texttt{cc12b5e1b600} & -- \\
HE/156 & 128 & 3/3 & 3/3 & \texttt{b8060d33b72c} & -- \\
HE/134 & 128 & 3/3 & 0/3 & \texttt{fc12c8453de5} & 88: 914 $\rightarrow$ 1946 \\
HE/42 & 128 & 3/3 & 0/3 & \texttt{c737973a39a6} & 73: 432 $\rightarrow$ 1817 \\
HE/122 & 128 & 3/3 & 3/3 & \texttt{ed3172303704} & -- \\
HE/114 & 128 & 3/3 & 3/3 & \texttt{16037f527b2f} & -- \\
HE/93 & 128 & 3/3 & 3/3 & \texttt{4b69a84d3675} & -- \\
HE/115 & 128 & 3/3 & 3/3 & \texttt{43fbb9bd3da9} & -- \\
HE/151 & 128 & 3/3 & 3/3 & \texttt{f371e4f2b94e} & -- \\
HE/63 & 128 & 3/3 & 3/3 & \texttt{b3141572ec95} & -- \\
HE/125 & 128 & 3/3 & 3/3 & \texttt{689af3f02f2d} & -- \\
HE/133 & 128 & 3/3 & 3/3 & \texttt{c1ed3bf83bde} & -- \\
HE/43 & 128 & 3/3 & 3/3 & \texttt{2e3776c7ff02} & -- \\
HE/139 & 128 & 3/3 & 3/3 & \texttt{23a77663974d} & -- \\
HE/77 & 128 & 3/3 & 3/3 & \texttt{dea569b38784} & -- \\
HE/35 & 128 & 3/3 & 3/3 & \texttt{93498855fe64} & -- \\
HE/27 & 128 & 3/3 & 3/3 & \texttt{f0f161a1b9e3} & -- \\
HE/34 & 128 & 3/3 & 3/3 & \texttt{813cb3180ab5} & -- \\
HE/145 & 128 & 3/3 & 0/3 & \texttt{31246ffe07b9} & 103: 1184 $\rightarrow$ 686 \\
HE/82 & 128 & 3/3 & 3/3 & \texttt{a356c9235624} & -- \\
HE/154 & 128 & 3/3 & 3/3 & \texttt{0743671d064c} & -- \\
HE/78 & 128 & 3/3 & 3/3 & \texttt{a685f47dbab9} & -- \\
HE/96 & 128 & 3/3 & 3/3 & \texttt{3824c19fc287} & -- \\
HE/138 & 128 & 3/3 & 0/3 & \texttt{9e54992a7160} & 93: 14374 $\rightarrow$ 9064 \\
\bottomrule
\end{tabular}
\end{table*}

The separate anonymous sequence supplement contains complete reference, predictor, and live generated-ID arrays for the 72 reserved runs, with a deterministic equality and divergence checker. It contains no identifying machine or repository metadata. Sanitizing the export changes its bytes, so it cannot reproduce the original prediction-file digest or independently prove the reported temporal freeze; the PDF's digest and chronology have those distinct stated roles.
\section{Follow-up Campaign Identities}
\label{app:followupconfigs}
% Evidence: *-PROTOCOL; manuscript evidence ledger.
\begin{table*}[t]
\caption{Independently frozen case studies; not a matched hardware comparison. All final target models are Qwen3-8B Q8\_0. Greedy execution uses campaign-matched target-only references; historical M4/RTX/B570 use re-tokenized text, while 358H uses direct IDs.}
\label{tab:platforms}
\centering
\begin{tabular}{>{\raggedright\arraybackslash}p{1.4in}>{\raggedright\arraybackslash}p{1.65in}>{\raggedright\arraybackslash}p{1.35in}>{\raggedright\arraybackslash}p{1.5in}}
\toprule
System & Runtime / target backend & Draft model / placement & Evidence scope \\
\midrule
MacBook Air M4 & llama.cpp b10809 / Metal & Qwen3-0.6B Q8\_0 / Metal & Representative 24Q and separate HumanEval30 \\
Core Ultra 9 275HX + RTX 5070 Ti Laptop & llama.cpp b10839 / CUDA & Qwen3-0.6B Q8\_0 / CPU or same GPU & Blocked fast path; text-path-agreeing controls \\
Core i5-13400 + Arc B570 & llama.cpp b10839 / Vulkan0 & Qwen3-0.6B Q8\_0 / B570 & Formal gate failure; later diagnostics \\
Core Ultra X7 358H + Arc B390 & llama.cpp b10839 / Vulkan0; OpenVINO draft & Qwen3-0.6B symmetric INT4 / CPU or NPU & Closed Phase 3; direct stateful decomposition \\
\bottomrule
\end{tabular}
\end{table*}

The follow-ups retain the Qwen3-8B/0.6B Q8\_0 pair and upstream llama.cpp b10839, but not one common frontend or reference. Table~\ref{tab:followupconfigs} identifies the new campaigns. Original RTX direct-ID qualification used Windows build 26200; new RTX H1, latency and H2 use 26300. Their driver version is 616.92, versus historical timing's 616.56; version equality does not authenticate original driver bytes. Recorded byte identities cover models, frontend, DLLs, compiler and observed driver modules on the experimental host; omitted assets cannot be rehashed by an offline reader.

\begin{table*}[t]
\caption{Campaign-specific follow-up configuration. RTX always denotes RTX 5070 Ti Laptop. Greedy target sampling uses seed 123, FA off, fresh contexts, and an inclusive EOG/128 committed-ID endpoint. All layers are offloaded: target 37/37 and draft 29/29. The original studies' references remain unchanged.}
\label{tab:followupconfigs}
\centering
\begin{tabular}{>{\raggedright\arraybackslash}p{1.15in}>{\raggedright\arraybackslash}p{2.15in}>{\raggedright\arraybackslash}p{2.45in}}
\toprule
Campaign & Frontend / execution & Reference or paired contrast \\
\midrule
B570 live fork & Instrumented native speculative example; Vulkan0; context 1024, batch 512 & Native UB1 versus state-restored UB128; identical fixed proposals at sampled block; UB128 continuation persists \\
RTX H1 discovery & Original CLI/CUDA DLLs; context 4096 & Target-only UB2 versus batched native live; 24/36 fails selection, reserved 0 \\
RTX H2 V2 & Separate P/R workers and original native CLI; shared numerical DLLs; context 4096, batch 2048 & H2 and R\_ONLY target UB2; original target-only reference UB1; native live target/draft UB2 \\
RTX matched cost & New Clang 21.1.0 timer frontend/core, original hashed CUDA backend DLLs; context 4096 & Same instrumented executable for target-only UB1 and serialized UB1; new 90/90 direct-ID gate \\
\bottomrule
\end{tabular}
\end{table*}

\paragraph{Separate selection and exposure.}
Original B570 H1's 24 reserved tasks and 133-task eligible pool remain as documented above. The live fork study excludes old direct-ID, old reserved and other recorded exposure before deterministic selection using its fixed study salt. The RTX H1 arm excludes 103 exposed or ambiguous tasks from 164, leaving 61; its reserved list is never executed after discovery failure. V2 re-audits accessible scientific files, archives and intervening study records: 114 tasks are excluded, leaving 50. Its selection uses the smallest SHA-256 values of a fixed V2 salt plus full task ID, without outcome selection. The actual 24-task V2 order is /160, /45, /149, /17, /6, /72, /33, /121, /130, /12, /4, /55, /69, /14, /61, /135, /100, /117, /22, /7, /23, /95, /88, /37. These tasks must not be pooled with original B570, the new fork set, the unexecuted H1 reserved list or latency's HumanEval30. Offline selection checks reconstruct retained inventories, not a new scan of unavailable remote directories.

\section{B570 Same-State Fork Endpoints}
The frozen checkpoint is the first native verification batch beginning at at least 64 committed IDs. It is one sampled block per eligible execution, not every generation position. All 72 reserved executions have eligible forks. The live draft's two proposals are recorded before either target decode. Source KV is serialized, restored into the alternate context, serialized again and compared with source-after bytes; consumed length, prefix, positions and physical groups are checked. Native UB1 produces three one-output physical groups, whereas alternate UB128 produces one three-output group at this sampled block.

Sampler identity uses the pinned recursive \texttt{common\_sampler\_clone}, copied chain/seed/accepted-token history and the same-partition R0 control. Six exposed engineering tasks generate 24 fresh R0 executions; original/instrumented primary IDs and proposals agree, and UB1-to-UB1 restored continuation and scores agree. Opaque sampler serialization is unavailable, so these are source-audited clone and empirical controls, not a byte proof for omitted external sampler binaries.

\begin{table}[t]
\caption{B570 reserved fork results: 24 new prompts $\times$ three. Events and repetitions are correlated. Primary immediate decisions remain unchanged; later output differences are a separate secondary endpoint.}
\label{tab:forkendpoints}
\centering
\begin{tabular}{lrr}
\toprule
Endpoint difference & Runs & Prompts \\
\midrule
Raw score records & \FUForkScores & 24 \\
Raw top-1 IDs & 0 & 0 \\
Immediate decisions & \FUForkImmediate & 0 \\
Full continuation IDs & \FUForkLater & \FUForkLaterPrompts \\
\bottomrule
\end{tabular}
\end{table}

The primary comparison follows every reachable decision at a still-common proposal-conditioned prefix in the sampled block until rejection, bonus or endpoint: \FUForkEvents\ paired events. All three raw score rows are retained even when proposals are rejected. Score equality, raw argmax equality and actual accepted/corrected/bonus IDs are different tests. The exported paired records preserve them separately.

After the native continuation ends, the implementation restores the saved draft KV and cloned samplers, replaces target context with the already-verified alternate UB128 context, restores counters and committed prefix, then continues the native loop. Subsequent target blocks remain UB128, and subsequent draft proposals follow alternate accepted history. This continuation policy changes more than a single numerical block over the rest of generation. The 18/72 trajectory differences across six prompts occur after unchanged sampled immediate decisions; they cannot be attributed to an immediate split at that block or claimed as reference restoration. At the tested states, the policy changes scores and later trajectories, with no detected immediate decision effect. These results identify no unique operation, do not rule out effects elsewhere in generation and support no RTX or performance inference.

\section{RTX Independent Proposal/Replay Reproduction}
\paragraph{The earlier failed arm.}
H1 target-only UB2 matches 24/36 newly executed exposed batched trajectories. /64, /106, /119 and /128 each fail all three repetitions. The 36/36 gate is unmet, no predictor is selected, no hypothesis or prediction freeze is reached, and reserved runs are zero. H2 was not implemented in that frozen arm. V2 is a new study, not a retroactive execution of the earlier decision tree.

\paragraph{Program independence and shared computation.}
P loads only draft weights; R loads only target weights. H2 orchestrates these processes, first recording each proposal request and then sending it to R. P uses H2's already-produced feedback; it is not a feedback-free draft tape. R computes pending tokens, acceptance, correction, bonus and KV suffix removal. Fresh R\_ONLY reads one sealed tape request at a time and checks its own computed prefix and pending token. Future prediction/reference/live IDs are excluded from its computation inputs. The original native CLI produces its own proposals and has no tape import. Model-loader, compile-time role, allowed RPC fields and process/load records establish application/API isolation, not adversarial OS-wide file access. Original model loaders, samplers, llama.cpp/GGML/CUDA kernels are shared, so there is no independent numerical implementation.

\paragraph{Freeze chain and denominators.}
F0 seals source, effective configuration, input/exposure/selection and qualification before discovery. Discovery completes 108 executions (12 tasks $\times$ three repetitions $\times$ H2, R\_ONLY and native live), with both H2/live and H2/R\_ONLY 36/36. F1 seals the selected program before 216 reserved pre-live executions (reference/H2/R\_ONLY). F2 seals the 72 prediction rows, full tapes and pre-live package; remote acknowledgement and retrieval of the exact sealed bytes precede the first of 72 reserved native live runs. Thus formal four-condition execution totals 288, with 72 comparisons per named endpoint. Git blob/tree membership certifies preserved bytes, while server and host records provide the reported order; neither is a trusted per-process external clock. Sanitized anonymous ordinals cannot independently prove that original absolute chronology.

\begin{table}[t]
\caption{RTX V2 reserved equality. Every row uses 24 distinct prompts $\times$ three repetitions. Prediction equality and reference qualification have different comparators.}
\label{tab:rtxreplay}
\centering
\begin{tabular}{lrr}
\toprule
Comparison & Equal / 72 & Prompts / 24 \\
\midrule
H2 / native live & 72 & 24 \\
R\_ONLY / native live & 72 & 24 \\
H2 / R\_ONLY & 72 & 24 \\
UB1 / native live & 60 & 20 \\
\bottomrule
\end{tabular}
\end{table}

\begin{table}[t]
\caption{RTX V2 reference-divergent cases, stable in all three repetitions. Indices are zero-based generated-ID positions; H2 and R\_ONLY predict the native live value. Reference is this study's UB1 path.}
\label{tab:rtxdivergences}
\centering
\begin{tabular}{lrrr}
\toprule
HumanEval & Index & Reference ID & Live ID \\
\midrule
/17 & 4 & 1406 & 271 \\
/33 & 8 & 14937 & 4210 \\
/14 & 4 & 1406 & 271 \\
/22 & 0 & 1499 & 750 \\
\bottomrule
\end{tabular}
\end{table}

\paragraph{Actual endpoint distribution.}
H2, R\_ONLY and native live each contain 63 earlier-EOG runs on 21 prompts and nine cap runs on three prompts. The UB1 reference has 66 earlier-EOG runs on 22 prompts and six cap runs on two prompts. Each includes the first EOG ID and excludes prompt and uncommitted proposals. This differs from original B570 H1, where every reserved array reaches 128. Repetition stability holds within all four conditions; it is not reference equality. The reference gate remains blocked. There is no target-only UB2 competitor on these same V2 reserved prompts, so these matches do not establish proposal/history necessity.

\section{Fresh RTX Serialized Cost Estimator and Timer}
Thirty distinct prompts from the latency arm's frozen original HumanEval30 input list each receive three executions of target-only and serialized: 180 new runs, 90 direct-ID pairs. Conditions are randomly ordered within prompt/repetition with seed 2026093001; bootstrap seed is 2026093002.

\begin{samepage}
For condition medians $m_{ic}$ on prompt $i$, compute
\[
\begin{aligned}
r_i&=m_{i,\mathrm{target}}/m_{i,\mathrm{serialized}},\\
G&=\exp\!\left(\frac1{30}\sum_{i=1}^{30}\log r_i\right).
\end{aligned}
\]
\end{samepage}
The 10,000 bootstrap draws resample 30 prompts in original numeric task order using Python Random, with linearly interpolated sorted percentiles at .025/.975. This is neither median of paired ratios nor ratio of 180 total times; repeats are not bootstrap units. Independent full-precision replay $G=0.6458157742698285$ (the sealed original report stores $0.6458157742698286$, one ulp higher), CI [0.629802205379344, 0.6598843928244866], with 0/30 ratios above one. The independent floating-point replay agrees within $10^{-14}$; full-precision canonical and independently recomputed results remain separately recorded. The rounded manuscript value is \FUCost\ [\FUCostLo, \FUCostHi].

Both conditions use the same freshly built timer executable and pass their own 90/90 direct-ID gate. The native QPC interval begins after synchronized target/draft initial prefill and before per-request sampler creation/reset/seed and generation-state work. Draft inference, target validation, acceptance/correction/bonus, KV state, synchronization, proposal/scheduling and relevant final response/debug-log work through slot release are charged. Loading, initial prefill and teardown are excluded. Persistent engine allocation and one-time draft sampler precede prefill; request sampler work is inside the timer. Endpoint records print after measured end. A second chrono interval cross-checks the native endpoints, but both APIs share Windows' clock oscillator.

The frozen empty-hook probe checks clock-call overhead at its prescribed threshold; it does not show that total committed-ID/debug logging costs less than 1\%. This is instrumented CLI cost, not deployment throughput. Thermal entry and randomized adjacent scheduling reduce drift without locking all clock/session variation. No energy result is estimated. The earlier RTX timing ratio 0.650, Figure 3 proxy cells and qualification-only gates are unchanged. The study has no qualified batched timing endpoint, so it does not measure a full batched-to-serialized speed tradeoff.

\subsection{Proposal inputs and local decisions}
\label{app:h2audit}
B570 H1 computes a target-only UB1 sequence with no proposals. RTX H2 instead constructs a feedback-dependent tape using draft-only P and target-only R; both share llama.cpp/GGML/CUDA. The sealed implementation uses this flow:
\begin{quote}\ttfamily
P(past feedback) -> proposals\par
R(own state, proposals) -> decisions\par
seal tape; fresh R\_ONLY -> own decisions
\end{quote}
R initialization accepts only the formatted prompt; its verification allowlist is \texttt{op, round, position, pending, proposal\_ids}. Round, position and pending must match R's locally maintained round, KV position and last sampled committed ID. Tape task/repetition, prefix count/hash and request dependencies are checked against the replay's own accumulated prefix. They are assertions of past state, not supplied future answers. Each tape row is consumed only after the previous verification completes; truncation, hash-chain mismatch, state mismatch, extra fields or a suffix after the endpoint aborts execution. R samples target logits locally, accepts matching proposals, emits correction at rejection or bonus after full acceptance, and removes the invalid suffix. P updates from its own verified batch and R's already-produced retain count. Live CLI generates its own proposals and imports no tape.

For reserved HumanEval/17 repetition 1, round 0 requests proposals [19496,1159] at position 195 with locally pending ID 1499. R samples [19496,1159,1759], accepts two proposals and commits bonus ID 1759 at generated index 3. Round 6 requests [4210,2505] at position 213, pending ID 262 and prefix length 19; R samples ID 8388, rejects the first proposal, retains one slot and commits ID 8388 at index 19. Both slices equal the final H2/R\_ONLY/live arrays. Source allowlists and sampler calls support the input separation; array equality alone does not prove non-leakage or independent numerical implementation.

The retained engineering capsule contains five modified-file witnesses (changed committed ID, deleted event, task mismatch, tape-round/hash mismatch, future-answer field) and a premature-live ordering witness, plus constructed acceptance/bonus and proposal-width fixtures. These are existing records, not new model tests. Offline rejection checks strengthen file/protocol review; they do not establish adversarial OS isolation. Process/model-loader records limit independence to application/API orchestration.

\section{Reviewer Evidence and Offline Recalculation}
\label{app:evidence}
\begin{table*}[t]
\caption{Evidence capabilities and exact reviewer entries. The source ZIP reconstructs documents; evidence packages recompute retained observations without models. Derived hashes/ordinals cannot certify original absolute chronology or omitted binaries.}
\label{tab:evidenceentries}
\centering
\begin{tabular}{@{}p{1.55in}p{2.80in}p{1.9in}@{}}
\toprule
Object & Exact package / entry & Capability and limit \\
\midrule
Document sources & \nolinkurl{overleaf-convergence-v0.8.6.zip}; \texttt{latexmk -pdf main.tex}, then \texttt{appendix.tex} & TeX dependency/hash audit and fresh build; no science execution \\
Compact studies & \nolinkurl{anonymous-evidence-index-v0.8.6.zip}; \texttt{python3 -B verify.py} & Complete ID arrays, scores/times and post-study inputs/statistics; no absolute time authentication \\
RTX full anonymous capsule & \nolinkurl{anonymous-a5919b59342ecc68bd5f57d3c2b192edf62933b2645924381c3765dd50beb369.tar.gz}; extract and run \nolinkurl{frozen_source/scripts/platform/rtx_independent_proposal_replay_v2/offline.py} with extracted root argument & Full proposals, tapes, events and code; shared numerical routines; sanitized ordinal ordering \\
Original freeze objects & Each study's source/policy, prediction/tape, acknowledgement and first-live records; compact \texttt{INDEX.json} distinguishes content roles & Retained bytes and reported host/server ordering; Git/hash identity alone is not trusted time \\
\bottomrule
\end{tabular}
\end{table*}
The H2 capsule's SHA256 is its full filename digest above, and its length is 60,152,666 bytes. \texttt{INDEX.json} gives separate-capsule identity and study-specific source/policy $\to$ predictions/tapes $\to$ acknowledgement $\to$ first-live $\to$ analysis roles. Original H1 has separate hypothesis and 72-array freezes; B570 forks freeze sampled policy before reserved execution; RTX H1 stops at discovery; RTX V2 F0/F1/F2 seal configuration, selected program, then predictions/tapes. RTX cost freezes its new timer and protocol before qualification/formal measurement; M4 policy recovery is documented in Section H. New descriptive analyses and exports follow known outcomes and claim no pre-live freeze.

Original B570 arrays are compared element by element with sealed originals; the exact earlier sequence ZIP is unavailable and is not replaced. The compact index adds post-study analysis inputs and a standard-library calculator while retaining existing arrays. Original host paths, identities, Git authors and absolute timestamps stay outside reviewer exports; internal reviews are excluded. Content digests may permit linkage and do not guarantee unlinkability. Retained records support reported original ordering; sanitized ordinal mapping cannot authenticate it independently. Omitted weights/runtime/compiler/driver bytes and full external assets cannot be rehashed by the offline reader. No functional HumanEval pass@$k$ or online upload is implied.

\section{Post-Study Existing-Record Analyses}
\label{app:derived}
The major outcomes were known before these v0.8.6 calculations. Before recomputation, a written plan fixes the complete retained inputs, formulas, ordering and treatment of missing records. These are secondary descriptive or sensitivity analyses; no original endpoint, primary CI or bootstrap stream is replaced. Anonymous \texttt{derived/inputs.json}, \texttt{derived/results.json} and \texttt{analysis.py} allow standard-library replay.

\subsection{B570 raw logits, margins and continuation}
Each reserved run retains three raw target-logit rows for a common proposal-conditioned prefix, whether or not all proposals are reached by actual acceptance. Top-one/top-two are ranked over the full vocabulary before sampler transformations; margin is their score difference. A fixed proposal's raw logit is also retained (absent for the bonus row). These logits are neither normalized probabilities nor sampler-adjusted margins. Numerical fields change in all 72 runs; timestamps/metadata are excluded. The 216 rows and 195 reachable paired decisions are correlated, as are three repeats per prompt.
\begin{table}[t]
\caption{B570 post-study raw-logit descriptors over all 216 paired rows. $\Delta$ = UB128 minus UB1; margins use retained top-two scores/IDs. Units are raw logits.}
\label{tab:scorederived}
\centering
\begin{tabular}{lrrr}
\toprule
Statistic & Min & Median & Max \\
\midrule
$\Delta$ top-one & -0.4275 & -0.0443 & 2.3503 \\
$|\Delta|$ top-one & 0.0019 & 0.0695 & 2.3503 \\
UB1 margin & 0.2291 & 17.5021 & 33.6660 \\
UB128 margin & 0.2003 & 17.5128 & 33.5351 \\
\bottomrule
\end{tabular}
\end{table}
Top-two ID changes occur in three rows; raw top-one IDs never change. Absolute top-one delta quartiles are 0.0388 and 0.1327, with p95 0.3437. Positive minimum raw margins coexist with unchanged immediate decisions; these summaries do not prove a whole-generation null effect or locate a numerical operation.

The new cohort is HumanEval/1, /8, /13, /21, /29, /36, /41, /51, /71, /74, /75, /79, /83, /84, /85, /92, /94, /110, /124, /131, /152, /153, /155, /163. Every primary and alternate endpoint reaches cap 128, with no earlier EOG. Source-before, restored and source-after serialized KV content hashes/lengths match at all 72 forks; consumed lengths and positions match. Recursive sampler cloning is source-audited, with matching recorded chain/seed/history. Six same-partition UB1-to-UB1 R0 restores have equal scores and full continuations; the complete R0 set has 24 existing executions including original/native equivalence. Sampler opaque-state bytes are not separately serialized. These controls support the sampled contrast, not identity of unavailable external binaries.

\begin{table}[t]
\caption{All six B570 continuation-sensitive prompts. Zero-based generated positions, distance = first difference minus fork index. Every row is stable across all three repetitions, with 128/128-ID capped endpoints. Length-only divergence would occur at the shorter array's length; none occurs here.}
\label{tab:forkpositions}
\centering
\begin{tabular}{lrrr}
\toprule
HumanEval & Fork & First difference & Distance \\
\midrule
/29 & 65 & 114 & 49 \\
/75 & 65 & 80 & 15 \\
/84 & 65 & 115 & 50 \\
/110 & 65 & 84 & 19 \\
/131 & 64 & 74 & 10 \\
/163 & 66 & 97 & 31 \\
\bottomrule
\end{tabular}
\end{table}
All sampled immediate decisions agree. Later UB128 target blocks persist and proposals evolve with the alternate accepted history. These secondary differences concern another cohort; they cannot explain old H1's four divergences or a single-block restoration.

\subsection{M4 opportunity between two fixed modes}
From all 540 raw held-out observations, the three-condition medians exactly reproduce the original 60 prompt rows. Target-only is faster than Always on 2/60 prompts, both Spec-Bench; there are no equal Target/Always medians. Always/oracle GM is 1.0006919671$\times$, where oracle is the per-prompt lower Target or Always median. It is limited hindsight choice between two measured fixed modes, not an arbitrary switching-policy lower bound.
\begin{table}[t]
\caption{The complete M4 target-faster subset, from original medians. Ratios are relative latencies ($\times$); lower is better. All six Guard runs fall back at round 8 but each Guard median remains above target-only. No transition component is inferred.}
\label{tab:m4opportunity}
\centering
\begin{tabular}{lrrr}
\toprule
Task & Target (s) & Always/Target & Guard/Target \\
\midrule
/559 & 11.9085 & 1.0270 & 1.0215 \\
/95 & 10.9203 & 1.0150 & 1.0029 \\
\bottomrule
\end{tabular}
\end{table}
Guard/Always paired prompt ratios have min/median/max 0.9706/1.0022/1.3297 and quartiles 0.9985/1.0260. Family GMs are HumanEval 1.0099 and Spec-Bench 1.0446; these 30/30 families are descriptive, not stable small-stratum estimates.

A separate sensitivity stream uses Python Random seed 2026100101, 10,000 paired draws, sampling each family's 30 prompts with replacement in HumanEval then Spec-Bench order. Each draw computes GM directly; linear .025/.975 percentiles give Guard/Always [1.0135, 1.0423] and Always/oracle [1.0000, 1.0018]. Guard remains slower; original pooled primary and standalone secondary intervals stay unchanged. No CI is formed by dividing interval endpoints.

\subsection{358H conditional zero-draft accounting}
For each path, sum matched baseline time $B$, complete target work $T$, draft $D_t$, residual $X_t$, committed IDs $N$ and rounds $L$ over all 90 processed observations. $T_b=B/N$, $C=N/L$, and $R_{\mathrm{target}}=T/(LT_b)$. At the same progress, target work and baseline, with nonnegative draft/residual,
\[
S=\frac{C}{D_{\mathrm{real}}+R_{\mathrm{target}}+X}
\le\frac{C}{R_{\mathrm{target}}}=\frac{B}{T}.
\]
Both paths retain 11,520 committed IDs; rounds are 4,677 CPU and 4,686 NPU. Full-precision processed times yield bounds CPU 0.8274529834 and NPU 0.8351858524$\times$. Even removing both draft and residual costs cannot reach parity under these fixed conditions. This is arithmetic total-time accounting, not a bound on prompt-GM speedup, a measured optimization, separable replay saving, persistence gain or a bound after changing target execution/proposals.

\subsection{RTX absolute instrumented latency}
All 180 original condition times are milliseconds; medians retain three repetitions per prompt and the original 30-prompt pairing. The following distributions summarize 30 condition medians, not 180 independent prompts.
\begin{table}[t]
\caption{RTX post-study absolute prompt-median times (ms). The original target/serialized GM and CI remain unchanged; the ratio of these distribution medians is not that GM.}
\label{tab:rtxabsolutetime}
\centering
\begin{tabular}{lrrrr}
\toprule
Condition & Min & Q25 & Median & Q75 \\
\midrule
Target & 295.2 & 1186.9 & 1907.9 & 2730.3 \\
Serialized & 560.3 & 1821.1 & 2923.3 & 4053.6 \\
\bottomrule
\end{tabular}
\end{table}
Maxima are 2939.7 and 4490.6 ms. The JSON output also reports all 90 individual times per condition separately. Independent primary replay is 0.6458157742698285; the original report's one-ulp 0.6458157742698286 and original CI/seed remain recorded. The post-prefill clock boundary and measured hook probe in Section N apply; 0.6 $\mu$s p99 is about 0.0012\% of the shortest 49.9572-ms R0 interval. That is a clock-call probe, not all logging overhead. The complete native-log files have median line counts 3516.5 target and 4353.5 serialized, and median byte lengths 266908 and 339335. Committed-ID event counts match for all 90 pairs. File-wide counts include initialization outside the timer; they do not measure charged logging cost or imply cancellation. No production-speed estimate is made.

\end{document}
